\documentclass[aspectratio=169,11pt]{beamer}

\usetheme{Madrid}
\usecolortheme{default}
\usefonttheme{professionalfonts}
\setbeamertemplate{navigation symbols}{}
\setbeamertemplate{footline}[frame number]

\usepackage{amsmath, amssymb, booktabs}

\title{Recruiting, Congestion, Timing, And Unraveling}
\subtitle{Process Frictions Inside Decentralized Hiring}
\author{Special Topic 7: Labor Market Design, Contracting, and Mechanism Design -- Week 2}
\date{}

\begin{document}

\begin{frame}
  \titlepage
\end{frame}

\begin{frame}{Central question}
\begin{itemize}
    \item Why do decentralized hiring markets often move too early, too fast, or too chaotically?
    \item The week opens the black box between vacancy posting and final hire.
    \item The labor objects are applications, interviews, wages, standards, offers, deadlines, acceptances, vacancy yields, and match quality.
    \item The research standard is translation: process friction $\rightarrow$ observed margin $\rightarrow$ data $\rightarrow$ counterfactual $\rightarrow$ welfare.
\end{itemize}
\end{frame}

\begin{frame}{Recruiting frictions versus fundamentals}
\begin{itemize}
    \item Fundamentals: worker skill, firm productivity, job amenities, location preferences, wage capacity.
    \item Recruiting process: posting design, wage disclosure, application costs, interview capacity, screening, offer timing, and deadlines.
    \item A bad final match can reflect weak fundamentals, but it can also reflect a poorly organized recruiting process.
    \item Key diagnostic: what margin failed before the match was observed?
    \item The design question is not ``who is good?'' only; it is also ``who gets seen, interviewed, compared, and allowed to wait?''
\end{itemize}
\end{frame}

\begin{frame}{Congestion, vacancy yields, and within-vacancy margins}
\small
\begin{tabular}{p{0.25\linewidth} p{0.34\linewidth} p{0.29\linewidth}}
\toprule
Margin & What changes inside a vacancy & Empirical signal \\
\midrule
Applicant pool & volume and composition & applications per posting, fit measures \\
Interview capacity & scarce evaluation slots & interview-to-application rate \\
Hiring standards & selectivity conditional on pool & offers per interview, unfilled jobs \\
Recruiting intensity & effort, wage, advertising, follow-up & vacancy duration, fill probability \\
Offer timing & deadlines and response windows & acceptance speed, yield loss \\
\bottomrule
\end{tabular}
\vspace{0.25cm}
\begin{itemize}
    \item Vacancy counts are not enough; vacancy-to-hire conversion is an equilibrium object.
\end{itemize}
\end{frame}

\begin{frame}{A simple recruiting funnel}
\begin{align*}
  A_v &= \text{applications to vacancy } v, \\
  I_v &= \text{interviews}, \\
  O_v &= \text{offers}, \\
  H_v &= \text{accepted hires}.
\end{align*}
\[
  \text{vacancy yield}_v = \frac{H_v}{A_v},
  \qquad
  \text{interview yield}_v = \frac{I_v}{A_v},
  \qquad
  \text{offer yield}_v = \frac{O_v}{I_v}.
\]
\begin{itemize}
    \item Different interventions move different ratios.
    \item A good empirical design says which ratio is the mechanism.
\end{itemize}
\end{frame}

\begin{frame}{Exploding offers, early contracting, and unraveling}
\begin{itemize}
    \item Firms make early offers because waiting risks losing candidates.
    \item Workers accept early because waiting risks having no offer later.
    \item Once some participants move early, others have incentives to move earlier too.
    \item The market loses thickness: choices are made before all participants and information arrive.
    \item Exploding offers shift risk to workers and can substitute deadline pressure for match quality.
\end{itemize}
\end{frame}

\begin{frame}{Professional labor-market evidence}
\small
\begin{tabular}{p{0.25\linewidth} p{0.34\linewidth} p{0.29\linewidth}}
\toprule
Market & Timing problem & Labor stake \\
\midrule
Gastroenterology fellowships & market with and without centralized match & mobility, wages, hiring practices \\
Clinical psychology & turnaround time and bottlenecks & interviews, early contracting, placement \\
Law firms and clerkships & informal norms and exploding offers & access, prestige, career starts \\
Residency transitions & early versus single match design & training slots and staffing \\
\bottomrule
\end{tabular}
\vspace{0.2cm}
\begin{itemize}
    \item These cases make timing rules visible as labor-market institutions.
\end{itemize}
\end{frame}

\begin{frame}{Modern vacancy, application, and job-post data}
\begin{itemize}
    \item Vacancy records: posting dates, wages, recruiter effort, duration, and closing outcomes.
    \item Job boards: applications, applicant characteristics, queue position, platform ranking, and posting content.
    \item Employer logs: screening, callbacks, interviews, offers, acceptances, and hires.
    \item Policy variation: wage transparency, application caps, preference signals, timing rules, or platform experiments.
    \item Linked outcomes: starting wages, retention, mobility, productivity, and worker welfare.
\end{itemize}
\end{frame}

\begin{frame}{Theory-to-empirical bridge}
\small
\begin{enumerate}
    \item Frictive object: congestion, wage opacity, interview bottleneck, unraveling, or low vacancy yield.
    \item Observed margin: applications, interview rate, offer timing, acceptance timing, duration, or hires.
    \item Data: vacancy, application, interview, offer, wage, and post-hire outcomes.
    \item Counterfactual: centralized match, timing rule, offer-hold period, wage disclosure, application cap, or queue design.
    \item Identification: regime change, staggered policy, experiment, platform rule, recruiter shock, or structural model.
\end{enumerate}
\vspace{0.1cm}
\begin{itemize}
    \item Keep equilibrium response visible: workers reallocate applications and firms adjust standards.
\end{itemize}
\end{frame}

\begin{frame}{Research Lab design}
\begin{itemize}
    \item Primary anchor: recruitment policies, job-filling rates, and matching efficiency.
    \item Challenge anchor: gastroenterology and timing/unraveling with and without centralized clearing.
    \item Reproduce: create vacancy-funnel facts about recruiting intensity, wage transparency, interview capacity, and yield.
    \item Diagnose: decompose low yield into applicant pool, screening, interview bottleneck, standards, and offer timing.
    \item Transfer: redesign the workflow for platforms, wage disclosure, clerkships, fellowships, or public-service recruiting.
\end{itemize}
\end{frame}

\begin{frame}{Takeaways}
\begin{itemize}
    \item Many matching failures are recruiting-process failures.
    \item Within-vacancy margins are central: applicant pools, interviews, standards, intensity, wage information, timing, and yield.
    \item Unraveling is a labor-market problem because it changes risk, information, access, mobility, and career starts.
    \item The frontier contribution is empirical: observe the process, name the margin, and evaluate a feasible design counterfactual.
\end{itemize}
\end{frame}

\end{document}
